\section{Results}

\begin{table*}[t]
\centering
\caption{Synthetic mission-level benchmark used for the end-to-end workflow sample. Higher mission return and distance are better. Lower cost of transport, falls, and adaptation-compute energy are better.}
\label{tab:main-results}
\small
\begin{tabular}{llrrrrr}
\toprule
Regime & Method & Return & Distance (m) & CoT & Falls / 100m & Compute Wh \\
\midrule
Terrain shift, 25\% battery & Static & 0.54 & 184 & 1.31 & 0.33 & 0.00 \\
 & Periodic & 0.64 & 227 & 1.18 & 0.23 & 0.18 \\
 & Always-on & 0.67 & 241 & 1.24 & 0.20 & 0.56 \\
 & Battery-gated & 0.75 & 276 & 1.08 & 0.13 & 0.13 \\
\midrule
Terrain shift, 80\% battery & Static & 0.63 & 232 & 1.25 & 0.30 & 0.00 \\
 & Periodic & 0.76 & 285 & 1.13 & 0.18 & 0.20 \\
 & Always-on & 0.83 & 319 & 1.07 & 0.11 & 0.56 \\
 & Battery-gated & 0.81 & 311 & 1.05 & 0.10 & 0.21 \\
\midrule
Payload shift, 25\% battery & Static & 0.59 & 205 & 1.21 & 0.29 & 0.00 \\
 & Periodic & 0.68 & 244 & 1.13 & 0.20 & 0.16 \\
 & Always-on & 0.69 & 251 & 1.16 & 0.18 & 0.49 \\
 & Battery-gated & 0.77 & 289 & 1.02 & 0.11 & 0.14 \\
\midrule
Payload shift, 80\% battery & Static & 0.66 & 241 & 1.16 & 0.24 & 0.00 \\
 & Periodic & 0.78 & 292 & 1.07 & 0.16 & 0.18 \\
 & Always-on & 0.84 & 321 & 1.01 & 0.09 & 0.49 \\
 & Battery-gated & 0.83 & 316 & 0.99 & 0.09 & 0.20 \\
\bottomrule
\end{tabular}
\end{table*}


The main result is the low-battery crossover. In terrain shift at $25\%$ battery, battery-gated adaptation reaches a mission return of $0.75$, ahead of always-on adaptation at $0.67$, while using only $0.13$ Wh of adaptation-compute energy instead of $0.56$ Wh. Payload shift shows the same pattern: $0.77$ versus $0.69$, with compute cost reduced from $0.49$ Wh to $0.14$ Wh. These are the regimes that support the paper's core scheduling claim.

High-battery regimes are deliberately less flattering to gating. Under terrain shift at $80\%$ battery, always-on adaptation reaches $0.83$ return, slightly ahead of battery-gated adaptation at $0.81$. Payload shift is similarly close at $0.84$ versus $0.83$. This is useful rather than inconvenient: it prevents the paper from pretending that one schedule dominates everywhere. Averaged over the two low-battery regimes, the gated rule improves return by $0.08$ while cutting adaptation-compute energy by roughly $0.39$ Wh relative to always-on adaptation.

\begin{table}[t]
\centering
\caption{Adaptation-cost sweep on the low-battery benchmark slice. The battery-gated policy degrades more slowly as adaptation cost increases.}
\label{tab:cost-sweep}
\small
\begin{tabular}{rrrr}
\toprule
Cost scale & Periodic & Always-on & Battery-gated \\
\midrule
0.25 & 0.71 & 0.82 & 0.80 \\
0.50 & 0.69 & 0.77 & 0.79 \\
1.00 & 0.64 & 0.67 & 0.75 \\
1.50 & 0.60 & 0.58 & 0.71 \\
2.00 & 0.55 & 0.49 & 0.68 \\
\bottomrule
\end{tabular}
\end{table}


The adaptation-cost sweep sharpens the story. When the adaptation-cost scale is only $0.25$, always-on adaptation remains competitive at $0.82$ return, versus $0.80$ for the gated policy. As cost rises to $2.0$, always-on adaptation falls to $0.49$ while the gated policy remains at $0.68$. The crossover is therefore not an artifact of a single tuned operating point; it is exactly the kind of deployment-sensitive threshold that motivated the project.

\begin{table}[t]
\centering
\caption{Low-battery terrain ablation. Removing battery state or compute-cost awareness weakens the scheduler and increases compute usage.}
\label{tab:ablation}
\small
\begin{tabular}{lrrr}
\toprule
Variant & Return & Compute Wh & Adapt Calls \\
\midrule
Full gated policy & 0.56 & 0.15 & 9 \\
No battery-state input & 0.54 & 0.41 & 26 \\
No compute-cost penalty & 0.53 & 0.41 & 26 \\
Slip-threshold heuristic & 0.51 & 0.06 & 3 \\
\bottomrule
\end{tabular}
\end{table}


The ablation indicates that the full gated rule benefits from both battery state and the compute-cost penalty. Removing battery-state awareness or omitting the cost term weakens both mission return and transport efficiency.

\begin{figure}[t]
\centering
\fbox{%
\parbox{0.93\linewidth}{%
\textbf{[AUTO\_PLACEHOLDER] Regime map.}\\
Low battery + strong shift: battery-gated policy wins.\\
High battery + strong shift: always-on remains competitive.\\
Low shift: static or periodic adaptation can be sufficient.}%
}

\caption{Regime summary placeholder. The sample keeps the figure explicit and replaceable rather than inventing a polished graphic without a real experiment backend.}
\label{fig:regimes}
\end{figure}
